\hypertarget{memory-types-and-pointers}{%
\chapter{MEMORY, TYPES, AND POINTERS}\label{memory-types-and-pointers}}

\hypertarget{advanced-pointers-memory}{%
\chapter{ADVANCED POINTERS \& MEMORY}\label{advanced-pointers-memory}}

\hypertarget{pointer-to-const-vs-const-pointer}{%
\section{1.1 Pointer to Const vs Const
Pointer}\label{pointer-to-const-vs-const-pointer}}

\begin{lstlisting}[language={C++}]
#include <iostream>
using namespace std;

int main() {
    int x = 5, y = 10;
    
    // Pointer to const - can't modify data
    const int* ptr1 = &x;
    // *ptr1 = 10;  // ERROR
    ptr1 = &y;      // OK - can change pointer
    
    // Const pointer - can't modify pointer
    int* const ptr2 = &x;
    *ptr2 = 10;     // OK - can change data
    // ptr2 = &y;   // ERROR
    
    // Const pointer to const - can't modify either
    const int* const ptr3 = &x;
    // *ptr3 = 10;  // ERROR
    // ptr3 = &y;   // ERROR
    
    return 0;
}
\end{lstlisting}

\hypertarget{void-pointers}{%
\section{1.2 Void Pointers}\label{void-pointers}}

\begin{lstlisting}[language={C++}]
#include <iostream>
using namespace std;

int main() {
    int x = 42;
    double y = 3.14;
    
    // Void pointer can point to any type
    void* ptr = &x;
    cout << *(int*)ptr << endl;  // 42
    
    ptr = &y;
    cout << *(double*)ptr << endl;  // 3.14
    
    // Generic function using void*
    void print_value(void* ptr, char type) {
        if (type == 'i') {
            cout << *(int*)ptr << endl;
        } else if (type == 'd') {
            cout << *(double*)ptr << endl;
        }
    }
    
    print_value(&x, 'i');  // 42
    print_value(&y, 'd');  // 3.14
    
    return 0;
}
\end{lstlisting}

\hypertarget{null-pointer-safety}{%
\section{1.3 Null Pointer Safety}\label{null-pointer-safety}}

\begin{lstlisting}[language={C++}]
#include <iostream>
using namespace std;

int main() {
    int* ptr = NULL;  // Set to null
    
    // Always check before dereferencing
    if (ptr != NULL) {
        cout << *ptr << endl;
    } else {
        cout << "Pointer is NULL" << endl;
    }
    
    // Safer approach
    ptr = new int(42);
    if (ptr) {
        cout << *ptr << endl;
        delete ptr;
        ptr = NULL;
    }
    
    return 0;
}
\end{lstlisting}

\hypertarget{memory-layout-alignment}{%
\section{1.4 Memory Layout \& Alignment}\label{memory-layout-alignment}}

\begin{lstlisting}[language={C++}]
#include <iostream>
using namespace std;

int main() {
    struct Data {
        char a;     // 1 byte
        int b;      // 4 bytes
        double c;   // 8 bytes
    };
    
    cout << "Size of Data: " << sizeof(Data) << endl;
    // Likely 16 or 24 (due to alignment padding)
    
    cout << "Size of char: " << sizeof(char) << endl;      // 1
    cout << "Size of int: " << sizeof(int) << endl;        // 4
    cout << "Size of double: " << sizeof(double) << endl;  // 8
    
    Data data;
    cout << "Address of a: " << (void*)&data.a << endl;
    cout << "Address of b: " << (void*)&data.b << endl;
    cout << "Address of c: " << (void*)&data.c << endl;
    
    return 0;
}
\end{lstlisting}

\begin{center}\rule{0.5\linewidth}{0.5pt}\end{center}

\hypertarget{advanced-arrays}{%
\chapter{ADVANCED ARRAYS}\label{advanced-arrays}}

\hypertarget{dynamic-arrays}{%
\section{4.1 Dynamic Arrays}\label{dynamic-arrays}}

\begin{lstlisting}[language={C++}]
#include <iostream>
using namespace std;

int main() {
    // 1D dynamic array
    int size = 5;
    int* arr = new int[size];
    
    for (int i = 0; i < size; i++) {
        arr[i] = i * 10;
    }
    
    for (int i = 0; i < size; i++) {
        cout << arr[i] << " ";
    }
    cout << endl;
    
    delete[] arr;
    arr = NULL;
    
    // 2D dynamic array
    int rows = 3, cols = 4;
    int** matrix = new int*[rows];
    for (int i = 0; i < rows; i++) {
        matrix[i] = new int[cols];
    }
    
    // Fill matrix
    for (int i = 0; i < rows; i++) {
        for (int j = 0; j < cols; j++) {
            matrix[i][j] = i * cols + j;
        }
    }
    
    // Delete matrix
    for (int i = 0; i < rows; i++) {
        delete[] matrix[i];
    }
    delete[] matrix;
    
    return 0;
}
\end{lstlisting}

\hypertarget{variable-length-arrays-non-standard}{%
\section{4.2 Variable Length Arrays
(Non-standard)}\label{variable-length-arrays-non-standard}}

\begin{lstlisting}[language={C++}]
#include <iostream>
using namespace std;

int main() {
    int size;
    cout << "Enter size: ";
    cin >> size;
    
    // VLA - not standard but supported by many compilers
    int arr[size];  // GCC extension
    
    for (int i = 0; i < size; i++) {
        arr[i] = i;
    }
    
    for (int i = 0; i < size; i++) {
        cout << arr[i] << " ";
    }
    cout << endl;
    
    return 0;
}
\end{lstlisting}

\hypertarget{array-bounds-safety}{%
\section{4.3 Array Bounds \& Safety}\label{array-bounds-safety}}

\begin{lstlisting}[language={C++}]
#include <iostream>
using namespace std;

int main() {
    int arr[5] = {10, 20, 30, 40, 50};
    
    // No bounds checking in C++
    cout << arr[0] << endl;   // 10 (OK)
    cout << arr[10] << endl;  // Undefined behavior!
    
    // Manual bounds checking
    int index = 5;
    if (index >= 0 && index < 5) {
        cout << arr[index] << endl;
    } else {
        cout << "Index out of bounds" << endl;
    }
    
    return 0;
}
\end{lstlisting}

\begin{center}\rule{0.5\linewidth}{0.5pt}\end{center}

\hypertarget{const-volatile}{%
\chapter{CONST \& VOLATILE}\label{const-volatile}}

\hypertarget{const-correctness}{%
\section{11.1 Const Correctness}\label{const-correctness}}

\begin{lstlisting}[language={C++}]
#include <iostream>
using namespace std;

int main() {
    int x = 10;
    
    // const variable
    const int constant = 5;
    // constant = 10;  // ERROR
    
    // pointer to const
    const int* ptr1 = &x;
    // *ptr1 = 20;  // ERROR
    ptr1 = &constant;  // OK
    
    // const pointer
    int* const ptr2 = &x;
    *ptr2 = 20;  // OK
    // ptr2 = &constant;  // ERROR
    
    // const reference
    const int& ref = x;
    // ref = 20;  // ERROR
    
    cout << x << endl;
    cout << *ptr1 << endl;
    cout << *ptr2 << endl;
    cout << ref << endl;
    
    return 0;
}
\end{lstlisting}

\hypertarget{volatile-keyword}{%
\section{11.2 Volatile Keyword}\label{volatile-keyword}}

\begin{lstlisting}[language={C++}]
#include <iostream>
using namespace std;

int main() {
    // volatile - tells compiler value may change unexpectedly
    volatile int sensor_reading = 0;  // From hardware
    
    // Compiler won't optimize away reads
    while (sensor_reading < 100) {
        // Check actual value each time, not cached
    }
    
    // Common use: hardware registers
    volatile int* hardware_register = (volatile int*)0x1000;
    
    // Each access reads from actual location
    int val1 = *hardware_register;
    int val2 = *hardware_register;
    
    // Without volatile, compiler might optimize one read
    
    return 0;
}
\end{lstlisting}

\begin{center}\rule{0.5\linewidth}{0.5pt}\end{center}

\hypertarget{type-casting}{%
\chapter{TYPE CASTING}\label{type-casting}}

\hypertarget{c-style-casting}{%
\section{8.1 C-Style Casting}\label{c-style-casting}}

\begin{lstlisting}[language={C++}]
#include <iostream>
#include <cmath>
using namespace std;

int main() {
    double d = 3.14;
    
    // C-style cast (avoid in modern C++)
    int i = (int)d;  // 3
    cout << i << endl;
    
    int x = 65;
    char c = (char)x;  // 'A'
    cout << c << endl;
    
    float f = (float)d;
    cout << f << endl;
    
    return 0;
}
\end{lstlisting}

\hypertarget{implicit-conversions}{%
\section{8.2 Implicit Conversions}\label{implicit-conversions}}

\begin{lstlisting}[language={C++}]
#include <iostream>
using namespace std;

int main() {
    // Implicit conversions
    int x = 5;
    double d = x;  // int to double (automatic)
    cout << d << endl;  // 5.0
    
    double d2 = 3.9;
    int y = d2;  // double to int (loses precision)
    cout << y << endl;  // 3
    
    // Char arithmetic
    char c = 'A';
    int code = c;  // char to int (gets ASCII)
    cout << code << endl;  // 65
    
    return 0;
}
\end{lstlisting}

\begin{center}\rule{0.5\linewidth}{0.5pt}\end{center}

\hypertarget{enumeration-unions}{%
\chapter{ENUMERATION \& UNIONS}\label{enumeration-unions}}

\hypertarget{enumerations}{%
\section{10.1 Enumerations}\label{enumerations}}

\begin{lstlisting}[language={C++}]
#include <iostream>
using namespace std;

// Enum definition
enum Color { RED, GREEN, BLUE };

enum Direction {
    NORTH = 0,
    EAST = 1,
    SOUTH = 2,
    WEST = 3
};

int main() {
    Color c = RED;
    cout << c << endl;  // 0
    
    Color colors[3] = {RED, GREEN, BLUE};
    
    // Switching on enum
    switch (c) {
        case RED:
            cout << "Red" << endl;
            break;
        case GREEN:
            cout << "Green" << endl;
            break;
        case BLUE:
            cout << "Blue" << endl;
            break;
    }
    
    // Iterate through enum values
    for (int dir = NORTH; dir <= WEST; dir++) {
        cout << "Direction: " << dir << endl;
    }
    
    return 0;
}
\end{lstlisting}

\hypertarget{unions}{%
\section{10.2 Unions}\label{unions}}

\begin{lstlisting}[language={C++}]
#include <iostream>
using namespace std;

// Union - all members share same memory
union Data {
    int i;
    float f;
    char c;
};

int main() {
    Data data;
    
    cout << "Size of Data: " << sizeof(data) << endl;  // 4 (size of largest member)
    
    data.i = 10;
    cout << "data.i: " << data.i << endl;     // 10
    cout << "data.f: " << data.f << endl;     // Garbage (overwrites data.i)
    
    data.f = 3.14;
    cout << "data.i: " << data.i << endl;     // Garbage (overwrites by data.f)
    cout << "data.f: " << data.f << endl;     // 3.14
    
    // Union useful for memory-constrained systems
    union Variant {
        int int_val;
        double double_val;
        char char_val;
    };
    
    cout << "Size of Variant: " << sizeof(Variant) << endl;
    
    return 0;
}
\end{lstlisting}

\begin{center}\rule{0.5\linewidth}{0.5pt}\end{center}

\hypertarget{bitwise-operations}{%
\chapter{BITWISE OPERATIONS}\label{bitwise-operations}}

\hypertarget{bitwise-operators}{%
\section{6.1 Bitwise Operators}\label{bitwise-operators}}

\begin{lstlisting}[language={C++}]
#include <iostream>
using namespace std;

int main() {
    unsigned char a = 5;   // 0101
    unsigned char b = 3;   // 0011
    
    // AND
    cout << (a & b) << endl;  // 0001 = 1
    
    // OR
    cout << (a | b) << endl;  // 0111 = 7
    
    // XOR
    cout << (a ^ b) << endl;  // 0110 = 6
    
    // NOT (bitwise complement)
    cout << (~a) << endl;     // 1010 = 250 (for unsigned char)
    
    // Left shift
    cout << (a << 1) << endl; // 1010 = 10
    
    // Right shift
    cout << (b >> 1) << endl; // 0001 = 1
    
    return 0;
}
\end{lstlisting}

\hypertarget{bit-manipulation-techniques}{%
\section{6.2 Bit Manipulation
Techniques}\label{bit-manipulation-techniques}}

\begin{lstlisting}[language={C++}]
#include <iostream>
using namespace std;

int main() {
    unsigned int num = 5;  // 0101
    
    // Check if bit is set
    int bit_pos = 2;
    bool is_set = (num >> bit_pos) & 1;
    cout << "Bit " << bit_pos << " is: " << is_set << endl;
    
    // Set a bit
    num |= (1 << 1);  // Set bit 1
    cout << "After setting bit 1: " << num << endl;  // 7 (0111)
    
    // Clear a bit
    num &= ~(1 << 1);  // Clear bit 1
    cout << "After clearing bit 1: " << num << endl;  // 5 (0101)
    
    // Toggle a bit
    num ^= (1 << 0);  // Toggle bit 0
    cout << "After toggling bit 0: " << num << endl;  // 4 (0100)
    
    // Count set bits
    unsigned int count = 0;
    unsigned int temp = num;
    while (temp) {
        count += temp & 1;
        temp >>= 1;
    }
    cout << "Number of set bits: " << count << endl;
    
    return 0;
}
\end{lstlisting}

\begin{center}\rule{0.5\linewidth}{0.5pt}\end{center}
